\documentclass[11pt]{article}

\usepackage[T1]{fontenc}
\usepackage[utf8]{inputenc}
\usepackage{microtype}
\usepackage{geometry}
\usepackage{xcolor}
\usepackage{amsmath,amssymb,mathtools}
\usepackage{booktabs}
\usepackage{tabularx}
\usepackage{array}
\usepackage{enumitem}
\usepackage{hyperref}
\usepackage{titlesec}
\usepackage{fancyhdr}
\usepackage{setspace}
\usepackage[most]{tcolorbox}
\usepackage{caption}
\usepackage{newtxtext,newtxmath}

\geometry{
  margin=0.96in,
  headheight=16pt,
  headsep=18pt,
  footskip=26pt
}

\definecolor{Ink}{HTML}{0B1220}
\definecolor{Muted}{HTML}{475467}
\definecolor{Line}{HTML}{D0D5DD}
\definecolor{Soft}{HTML}{F8FAFC}
\definecolor{SoftAlt}{HTML}{F2F4F7}
\definecolor{Accent}{HTML}{0F766E}
\definecolor{AccentDeep}{HTML}{134E4A}
\definecolor{AccentWash}{HTML}{ECFDF3}

\hypersetup{
  colorlinks=true,
  linkcolor=AccentDeep,
  citecolor=AccentDeep,
  urlcolor=AccentDeep,
  pdftitle={A Methodology for Developmental Trajectory Analysis of Attention Heads},
  pdfauthor={Abderahmane Ainouche}
}

\setstretch{1.05}
\setlength{\parindent}{0pt}
\setlength{\parskip}{0.42em}
\setlist[itemize]{leftmargin=1.3em,itemsep=0.25em,topsep=0.4em}
\setlist[enumerate]{leftmargin=1.45em,itemsep=0.25em,topsep=0.4em}
\captionsetup{labelfont=bf,textfont=small}

\newcolumntype{Y}{>{\raggedright\arraybackslash}X}

\titleformat{\section}
  {\normalfont\sffamily\Large\bfseries\color{Ink}}
  {\thesection}
  {0.75em}
  {}
  [\vspace{0.2em}\titlerule]

\titleformat{\subsection}
  {\normalfont\sffamily\large\bfseries\color{Ink}}
  {\thesubsection}
  {0.65em}
  {}

\titleformat{\subsubsection}
  {\normalfont\sffamily\normalsize\bfseries\color{Ink}}
  {\thesubsubsection}
  {0.6em}
  {}

\pagestyle{fancy}
\fancyhf{}
\fancyhead[L]{\textcolor{Muted}{\sffamily\small Head Trajectories Methodology}}
\fancyhead[R]{\textcolor{Muted}{\sffamily\small \thepage}}
\renewcommand{\headrulewidth}{0.35pt}
\renewcommand{\headrule}{\hbox to\headwidth{\color{Line}\leaders\hrule height \headrulewidth\hfill}}

\tcbset{
  enhanced,
  boxrule=0.5pt,
  colframe=Line,
  arc=3pt,
  left=9pt,
  right=9pt,
  top=8pt,
  bottom=8pt
}

\newtcolorbox{abstractbox}{
  colback=Soft,
  colframe=Line
}

\newtcolorbox{accentbox}[1]{
  colback=AccentWash,
  colframe=Accent,
  title=\sffamily\bfseries #1,
  coltitle=AccentDeep,
  boxrule=0.7pt
}

\newtcolorbox{roadmapbox}{
  colback=SoftAlt,
  colframe=Line
}

\newcommand{\E}{\mathbb{E}}
\newcommand{\Var}{\mathrm{Var}}
\newcommand{\Cov}{\mathrm{Cov}}
\newcommand{\KL}{D_{\mathrm{KL}}}
\newcommand{\1}{\mathbf{1}}
\newcommand{\RR}{\mathbb{R}}
\newcommand{\pipelinearrow}{\;\rightarrow\;}

\begin{document}
\thispagestyle{empty}

\vspace*{0.4cm}
{\sffamily\color{AccentDeep}\small\bfseries METHODS NOTE}

\vspace{0.55cm}

{\sffamily\fontsize{25}{30}\selectfont\bfseries\color{Ink}
A Methodology for Developmental Trajectory\\[-0.05em]
Analysis of Attention Heads
\par}

\vspace{0.5cm}

{\sffamily\large\color{Muted}
Formalizing the measurement instrument, behavioral metrics, calibration null,
classification rule, and trajectory analyses used to study how attention heads
emerge during transformer training
\par}

\vspace{0.9cm}

\begin{tabularx}{\textwidth}{@{}p{0.16\textwidth}Y p{0.14\textwidth}Y@{}}
\textcolor{Muted}{\sffamily Author} & Abderahmane Ainouche &
\textcolor{Muted}{\sffamily Scope} & Theory, concepts, and methodology only \\
\textcolor{Muted}{\sffamily Academic} & \href{mailto:abderahmane.ainouche@ensia.edu.dz}{abderahmane.ainouche@ensia.edu.dz} &
\textcolor{Muted}{\sffamily Contact} & \href{mailto:abderahmane.ainouche.ai@gmail.com}{abderahmane.ainouche.ai@gmail.com} \\
\textcolor{Muted}{\sffamily Project} & Developmental Trajectories of Attention Heads &
\textcolor{Muted}{\sffamily Repository} & \url{https://github.com/abderahmane-ai/head-trajectories} \\
\end{tabularx}

\vspace{0.9cm}

\begin{abstractbox}
{\sffamily\bfseries\color{Ink} Abstract}

Most attention-head interpretability is static: a trained model is inspected and
heads are assigned functions after optimization has finished. This project asks
a developmental question instead: \emph{when do recognizable head behaviors
emerge, and how do they evolve over training?} The answer requires more than
post-hoc interpretation. It requires a fixed probe dataset, stable score
definitions, a principled baseline for non-random behavior, and a trajectory
representation that can be analyzed consistently across checkpoints.

This document formalizes that methodology. A transformer is checkpointed
throughout training, a fixed held-out probe set is run through every checkpoint,
each head is scored on five behaviors, thresholds are calibrated from a
causally valid random baseline, and heads are classified and tracked over time.
The output is not merely a snapshot of head types, but a family of calibrated
head trajectories from which onset, stability, layer stratification, and
transition structure can be studied. The goal of this paper is methodological
clarity: to specify exactly what is measured, how it is measured, and what
interpretive claims the pipeline can and cannot support.
\end{abstractbox}

\vspace{0.7cm}

\begin{roadmapbox}
{\sffamily\bfseries\color{Ink} Method at a glance}
\[
\text{train with checkpoints}
\pipelinearrow
\text{run fixed probes}
\pipelinearrow
\text{score each head}
\pipelinearrow
\text{calibrate thresholds}
\pipelinearrow
\text{classify heads}
\pipelinearrow
\text{analyze trajectories.}
\]
\end{roadmapbox}

\vspace{0.9cm}

\section{Overview}

The Head Trajectories project treats developmental interpretability as a
measurement problem. The central challenge is not simply defining interesting
head behaviors, but building an experimental pipeline in which changes over time
can be attributed to the model rather than to drift in the instrument.

The methodology is built around four commitments:

\begin{enumerate}
\item the probe dataset is fixed across checkpoints and runs;
\item behavioral scores are defined once and applied uniformly;
\item thresholds are calibrated relative to an architecture-matched random
baseline;
\item the main object of analysis is a trajectory, not a terminal label.
\end{enumerate}

\begin{accentbox}{Design principle}
The pipeline should make the statement ``this head became more sink-like at
step $t$'' meaningful in a calibrated sense. That requires fixed probes,
architecture-matched baselines, and operational definitions of onset and
specialization.
\end{accentbox}

\section{Experimental Objects and Notation}

Let $\theta^{(k)}$ denote model parameters at checkpoint index
$k \in \{1,\dots,K\}$. The model has $L$ layers, $H$ attention heads per layer,
sequence length $T$, and vocabulary size $V$.

The project studies the ordered family
\[
\theta^{(1)} \pipelinearrow \theta^{(2)} \pipelinearrow \cdots \pipelinearrow \theta^{(K)}.
\]

For checkpoint $k$, layer $\ell$, head $h$, and probe sequence $n$, let
\[
A_n^{(k,\ell,h)} \in [0,1]^{T \times T}
\]
denote the post-softmax causal attention matrix. Because the model is
decoder-only,
\[
A_n^{(k,\ell,h)}[t,j] = 0 \qquad \text{for } j > t,
\]
and each query row is stochastic:
\[
\sum_{j=0}^{T-1} A_n^{(k,\ell,h)}[t,j] = 1.
\]

\begin{table}[h]
\centering
\renewcommand{\arraystretch}{1.18}
\begin{tabularx}{\textwidth}{@{}p{0.18\textwidth}p{0.18\textwidth}Y@{}}
\toprule
\textbf{Symbol} & \textbf{Type} & \textbf{Meaning} \\
\midrule
$\theta^{(k)}$ & parameters & model state at checkpoint $k$ \\
$A_n^{(k,\ell,h)}$ & $T \times T$ matrix & attention map for head $(\ell,h)$ on sequence $n$ \\
$s^{(k,\ell,h)}$ & 5-vector & behavioral score vector for one head \\
$\tau$ & 5-vector & calibrated threshold vector \\
$y^{(\ell,h)}$ & length-$K$ sequence & label trajectory for one head across checkpoints \\
$f_{k,c}$ & scalar & global fraction of heads of type $c$ at checkpoint $k$ \\
\bottomrule
\end{tabularx}
\caption{Core objects used throughout the methodology.}
\end{table}

\section{Probe Dataset as a Fixed Instrument}

The probe dataset is held fixed across checkpoints, seeds, and analyses. This
immutability is essential. If the instrument changes, trajectory claims become
ambiguous.

The dataset is partitioned into three components.

\subsection{General probes}

General probes are ordinary held-out token sequences
\[
X^{\mathrm{gen}} = \{x^{\mathrm{gen}}_1,\dots,x^{\mathrm{gen}}_{N_g}\},
\qquad
x^{\mathrm{gen}}_n \in \{0,\dots,V-1\}^T.
\]
They support sink, previous-token, and semantic scoring.

\subsection{Induction probes}

Induction probes are held-out sequences modified to contain a repeated
subsequence. For each probe $n$, the dataset stores
\[
x_n^{\mathrm{ind}}, \qquad p_n^{(1)}, \qquad p_n^{(2)},
\]
where $p_n^{(1)}$ and $p_n^{(2)}$ index the first and second occurrences of the
repeated pattern.

\subsection{Positional probes}

Positional probes are content-distinct sequence pairs of equal length:
\[
\bigl(x_i^{(a)},x_i^{(b)}\bigr), \qquad i=1,\dots,N_p.
\]
These test whether a head's structure is primarily position-driven rather than
content-driven.

\begin{accentbox}{Why fixed probes matter}
The methodology is longitudinal. It is not enough to say that a probe dataset is
``reasonable.'' The same dataset must be reused so that score movement over time
reflects changing model behavior rather than changing measurement conditions.
\end{accentbox}

\section{Behavioral Metrics}

For a fixed checkpoint $k$, layer $\ell$, and head $h$, define the score vector
\[
s^{(k,\ell,h)} =
\bigl(s_{\mathrm{sink}}, s_{\mathrm{prev}}, s_{\mathrm{ind}},
s_{\mathrm{pos}}, s_{\mathrm{sem}}\bigr).
\]
Below, dependence on $(k,\ell,h)$ is suppressed for readability.

\subsection{Sink score}

The sink metric detects fixed-key anchoring. A head should score high only if
many queries route attention to the same absolute key position.

\[
s_{\mathrm{sink}}
=
\max_{j \in \{0,\dots,T-1\}}
\left(
\frac{1}{N_g}
\sum_{n=1}^{N_g}
\frac{\sum_{t=0}^{T-1} A_n^{\mathrm{gen}}[t,j]}{T-j}
\right).
\]

The denominator $(T-j)$ corrects for causal geometry: key $j$ is reachable from
exactly $(T-j)$ query rows.

\subsection{Previous-token score}

The previous-token score reads the mass placed on the subdiagonal:
\[
s_{\mathrm{prev}}
=
\frac{1}{N_g(T-1)}
\sum_{n=1}^{N_g}
\sum_{t=1}^{T-1}
A_n^{\mathrm{gen}}[t,t-1].
\]

\subsection{Induction score}

The induction score measures whether the head attends from the repeated pattern
to what followed its earlier occurrence:
\[
s_{\mathrm{ind}}
=
\frac{1}{N_i}
\sum_{n=1}^{N_i}
A_n^{\mathrm{ind}}\bigl[p_n^{(2)},\; p_n^{(1)} + 1\bigr].
\]

\subsection{Positional score}

The current production positional metric compares paired attention maps using a
clipped rowwise KL-divergence similarity. Let $A_i$ and $B_i$ denote the two
attention matrices for positional pair $i$. Define
\[
\widetilde{B}_i[t,j]
=
\frac{B_i[t,j]+\varepsilon}{\sum_{u=0}^{T-1}(B_i[t,u]+\varepsilon)},
\qquad
\varepsilon = 10^{-9},
\]
and
\[
\overline{D}_i
=
\frac{1}{T}
\sum_{t=0}^{T-1}
\KL\!\left(A_i[t,\cdot]\;\|\;\widetilde{B}_i[t,\cdot]\right).
\]
Then
\[
s_{\mathrm{pos}}^{(i)} = \max(0,\,1-\overline{D}_i),
\qquad
s_{\mathrm{pos}} = \frac{1}{N_p}\sum_{i=1}^{N_p} s_{\mathrm{pos}}^{(i)}.
\]

\subsection{Semantic score}

Let $E \in \RR^{V \times d}$ be the checkpoint embedding matrix and let
\[
\widehat{E}_v = \frac{E_v}{\lVert E_v \rVert_2}
\]
denote the row-normalized embedding for token $v$.

For general probe sequence
\[
x_n = (x_{n,0},\dots,x_{n,T-1}),
\]
and for query position $t \ge 4$, define
\[
a_{n,t} = \bigl(A_n^{\mathrm{gen}}[t,0],\dots,A_n^{\mathrm{gen}}[t,t]\bigr)
\]
and
\[
c_{n,t}
=
\bigl(
\cos(\widehat{E}_{x_{n,t}},\widehat{E}_{x_{n,0}}),\dots,
\cos(\widehat{E}_{x_{n,t}},\widehat{E}_{x_{n,t}})
\bigr).
\]

Before correlation is computed, three confounded positions are removed:
\[
j \in \{0,\; t-1,\; t\}.
\]
These correspond to the sink position, the previous-token position, and the
trivial identity position. Let $M_{n,t}$ be the resulting mask. Only positions
with at least six unmasked entries are retained.

The semantic score is the empirical mean Pearson correlation over valid
sequence-position pairs:
\[
s_{\mathrm{sem}}
=
\E_{(n,t)\in\mathcal{V}}
\left[
\frac{\Cov(a_{n,t}\odot M_{n,t},\,c_{n,t}\odot M_{n,t})}
{\sqrt{\Var(a_{n,t}\odot M_{n,t})}\sqrt{\Var(c_{n,t}\odot M_{n,t})}}
\right].
\]

\begin{accentbox}{Interpretive note}
The positional metric is intentionally treated as an operational measurement rule
rather than a perfect ontological definition of positionality. It is useful
enough for trajectory analysis, but alternative similarity measures remain a
valid direction for follow-up evaluation.
\end{accentbox}

\section{Calibration Against Random Baselines}

\subsection{Why calibration is necessary}

The five metrics live on different scales. A raw score of $0.08$ can be
uninformative for one metric and strong for another. Direct raw comparison is
therefore not scientifically meaningful.

The calibration problem is:
\begin{center}
\emph{How large must a score be before it stops looking like random behavior for
this architecture and this metric?}
\end{center}

\subsection{The null transformation}

For a given architecture, calibration:
\begin{enumerate}
\item initializes a random model,
\item runs the fixed probe dataset through it,
\item extracts all attention maps,
\item causally scrambles key positions within each valid query row,
\item recomputes all five metrics for every head.
\end{enumerate}

The choice of null is deliberate. Row shuffling is not sufficient for the sink
metric because sink scoring aggregates over query rows. The null must destroy
fixed-key anchoring itself. The adopted transformation therefore permutes
\emph{key positions within each valid causal row} and then renormalizes the row:
\[
A_n[t,0:t] \mapsto \Pi_{n,t}A_n[t,0:t],
\]
where $\Pi_{n,t}$ is a permutation matrix acting only on the valid causal
support of row $t$.

\begin{accentbox}{Null-model principle}
The calibration null should destroy the exact structure a metric is designed to
detect while preserving enough geometry to remain a valid attention object:
causal support and row-stochasticity are preserved, but interpretable structure
is broken.
\end{accentbox}

\subsection{Threshold construction}

Let $r_{u,m}$ denote the random-baseline score of head $u$ on metric $m$, where
$u$ indexes all heads in the architecture. Define
\[
\mu_m = \frac{1}{LH}\sum_{u=1}^{LH} r_{u,m},
\qquad
\sigma_m =
\sqrt{
\frac{1}{LH}
\sum_{u=1}^{LH}(r_{u,m}-\mu_m)^2
}.
\]
The single-seed threshold is
\[
\tau_m = \mu_m + 2\sigma_m.
\]

Calibration is repeated across several random seeds. The repository stores:
\begin{itemize}
\item per-seed threshold vectors,
\item per-seed metric means,
\item per-seed metric standard deviations,
\item whether any threshold is non-positive,
\item the mean threshold vector across seeds,
\item the threshold standard deviation across seeds.
\end{itemize}

\section{Thresholded Classification}

Let
\[
\tau = (\tau_1,\dots,\tau_5)
\]
be the raw calibrated threshold vector and
\[
s = (s_1,\dots,s_5)
\]
the observed score vector for a head.

\subsection{Safe threshold handling}

Threshold summaries are still stored for diagnostics and legacy compatibility,
but the primary classifier gate is now statistical inference against the pooled
empirical null distribution.

\subsection{Decision rule}

For score vector $s=(s_1,\dots,s_5)$ and pooled null samples, we compute
one-sided empirical p-values:
\[
p_m = \Pr_{\mathrm{null}}(r_m \ge s_m).
\]

Per head and checkpoint, BH-FDR is applied across the five metrics at level
$\alpha$ to define the active behavior set:
\[
\mathcal{A} = \{m \in \{1,\dots,5\} : p_m \text{ survives BH-FDR}\}.
\]

For dominant summarization, effect sizes are
\[
e_m = -\log_{10}(p_m).
\]
Let $e_{(1)} \ge e_{(2)}$ be the top two active effect sizes and
\[
\Delta = e_{(1)} - e_{(2)}.
\]
The dominant-summary label set is
\[
\mathcal{Y}
=
\{\texttt{WEAK},\texttt{AMBIGUOUS},\texttt{SINK},\texttt{PREV\_TOKEN},
\texttt{INDUCTION},\texttt{POSITIONAL},\texttt{SEMANTIC}\}.
\]

Classifier:
\[
\mathrm{label}(s)
=
\begin{cases}
\texttt{WEAK}, & \mathcal{A}=\varnothing, \\
\texttt{AMBIGUOUS}, & |\mathcal{A}|>1 \land \Delta < \delta, \\
\arg\max_{m \in \mathcal{A}} e_m, & \text{otherwise.}
\end{cases}
\]

This keeps detection (active set) and summarization (dominant label) separate.

\section{Trajectory-Level Analysis}

\subsection{Global curves}

For type $c \in \mathcal{Y}$, the fraction of heads of type $c$ at checkpoint
$k$ is
\[
f_{k,c}
=
\frac{1}{LH}
\sum_{\ell=1}^{L}
\sum_{h=1}^{H}
\1\{y_k^{(\ell,h)} = c\}.
\]

\subsection{Per-layer curves}

Layer-restricted fractions are
\[
f_{k,c}^{(\ell)}
=
\frac{1}{H}
\sum_{h=1}^{H}
\1\{y_k^{(\ell,h)} = c\}.
\]
These support layer-stratification analysis.

\subsection{Transitions and stability}

The number of type changes for a head is
\[
\Delta^{(\ell,h)}
=
\sum_{k=2}^{K}
\1\{y_k^{(\ell,h)} \neq y_{k-1}^{(\ell,h)}\}.
\]
This statistic quantifies stability and re-specialization.

\subsection{Onset}

Given threshold fraction $\alpha$, the onset of type $c$ is
\[
\mathrm{onset}_c(\alpha) = \min\{k : f_{k,c} \ge \alpha\}.
\]

The project typically uses $\alpha = 0.05$, but positional onset requires a
special interpretive split.

\begin{accentbox}{Architectural vs.\ learned positionality}
The classifier keeps a single \texttt{POSITIONAL} label. The distinction between
architectural and learned positional structure is made in analysis:
\begin{itemize}
\item \textbf{architectural positional onset} is the immediate step-0 structure
driven by RoPE;
\item \textbf{learned positional onset} is the first later crossing after the
step-0 architectural effect is excluded.
\end{itemize}
This keeps the classifier simple while preserving a cleaner interpretation of
developmental order.
\end{accentbox}

\section{What the Method Measures}

The methodology does \emph{not} claim to recover a head's true essence. It
measures behavioral specialization relative to a fixed probe set and a
calibrated random baseline.

Concretely, it measures:
\begin{itemize}
\item whether a head looks sink-like, previous-token-like, induction-like,
positional, or semantic on the probe set;
\item whether that behavior is large relative to the null distribution for the
same architecture;
\item how that classification changes over training.
\end{itemize}

It does \emph{not} guarantee:
\begin{itemize}
\item that the resulting label is exhaustive;
\item that the label is unique outside the probe regime;
\item that onset is anything more than an operational threshold crossing;
\item that the five metrics span all meaningful head behaviors.
\end{itemize}

\section{Limits and Extensions}

\subsection{Strengths}

\begin{itemize}
\item the probe set is fixed across time;
\item the null model is architecture-matched;
\item the sink metric now uses a calibration null that actually destroys
fixed-key anchoring;
\item structural confounds are explicitly removed from the semantic metric;
\item both raw and effective thresholds are recorded, so defensive fallbacks are
auditable.
\end{itemize}

\subsection{Limits}

\begin{itemize}
\item the current positional metric is still heuristic;
\item argmax classification compresses potentially mixed behaviors into one
label;
\item onset depends on an operational threshold fraction;
\item cross-seed agreement remains an empirical question, not a methodological
guarantee.
\end{itemize}

\subsection{Natural extensions}

\begin{itemize}
\item richer positional similarity measures,
\item multi-label or continuous head identities,
\item broader behavior libraries beyond the current five metrics,
\item trajectory comparisons across architectures and datasets.
\end{itemize}

\section{Conclusion}

The Head Trajectories project is fundamentally a measurement pipeline for
developmental interpretability. Models are checkpointed through training, a
fixed probe set is reused at every checkpoint, five behavior scores are
computed for every head, thresholds are calibrated with a causally valid random
baseline, and head identities are tracked as trajectories. The result is a
framework designed to answer a developmental question in a disciplined way: not
only what attention heads do after training, but how recognizable behaviors
appear, stabilize, and transform during optimization.

\vspace{0.8em}
\begin{abstractbox}
{\sffamily\bfseries\color{Ink} One-sentence summary}

The methodology measures attention-head development by scoring every head at
every checkpoint against a fixed probe dataset, calibrating behavioral
significance with a causally valid random baseline, and analyzing the resulting
label trajectories over training.
\end{abstractbox}

\bibliographystyle{plain}
\begin{thebibliography}{9}

\bibitem{vaswani2017attention}
Vaswani, A., Shazeer, N., Parmar, N., Uszkoreit, J., Jones, L., Gomez, A. N.,
Kaiser, \L., and Polosukhin, I. (2017).
\newblock Attention is all you need.
\newblock In \emph{Advances in Neural Information Processing Systems}.

\bibitem{elhage2021mathematical}
Elhage, N., Nanda, N., Olsson, C., Henighan, T., Joseph, N., Mann, B., Askell,
A., Bai, Y., Chen, A., Conerly, T., et al. (2021).
\newblock A mathematical framework for transformer circuits.
\newblock \emph{Transformer Circuits Thread}.

\bibitem{olsson2022context}
Olsson, C., Elhage, N., Nanda, N., Joseph, N., DasSarma, N., Henighan, T.,
Mann, B., Askell, A., Bai, Y., Chen, A., et al. (2022).
\newblock In-context learning and induction heads.
\newblock \emph{arXiv preprint arXiv:2209.11895}.

\bibitem{belinkov2017neural}
Belinkov, Y., Durrani, N., Dalvi, F., Sajjad, H., and Glass, J. (2017).
\newblock What do neural machine translation models learn about morphology?
\newblock In \emph{Proceedings of ACL}.

\bibitem{linzen2016assessing}
Linzen, T., Dupoux, E., and Goldberg, Y. (2016).
\newblock Assessing the ability of LSTMs to learn syntax-sensitive dependencies.
\newblock \emph{Transactions of the Association for Computational Linguistics}.

\end{thebibliography}

\end{document}
